\documentclass[11pt]{article}
\usepackage[margin=0.85in]{geometry}
\usepackage{amsmath,amssymb,booktabs,longtable,array,hyperref,microtype,xcolor}
\usepackage[T1]{fontenc}
\usepackage{lmodern}
\hypersetup{colorlinks=true,linkcolor=blue!45!black,urlcolor=blue!45!black}
\title{FinCompass User Guide}
\author{Rajeev Yadav, Ph.D.}
\date{Version 2.0}
\begin{document}
\maketitle
\tableofcontents
\newpage
\section{What FinCompass does}
FinCompass helps you research a security, understand its historical behavior, and---when an appropriate model exists---estimate the probability of a precisely defined future event. It is local-first: models, watch-state, research history, and most calculations remain on your computer.

FinCompass separates three things that should never be confused:
\begin{enumerate}
\item an evidence score describing the information currently observed;
\item a Forecast probability for a defined future event;
\item a Live probability that may adjust the validated Forecast only after additional evidence matures.
\end{enumerate}

A probability such as 62\% does not mean the stock will rise 62\%, that the stock is 62\% safe, or that the model is 62\% accurate.

\section{The shortest useful workflow}
A normal workflow is:
\begin{enumerate}
\item Enter a company or ticker.
\item Choose a forecast period such as 6 months, 1 year, 2 years, or 3 years.
\item Let FinCompass identify the market, benchmark, local-data state, and eligible model.
\item If data are old, choose \textbf{Update data}.
\item If a model is old, choose \textbf{Update model}.
\item Read the Forecast and its evidence label.
\item Start Live only when the model tier is eligible.
\end{enumerate}

The software should not require you to understand model IDs, feature contracts, calibration algorithms, or activation pointers to complete this workflow.


\section{A new ticker should feel simple}
For an ordinary user, entering a new supported ticker should not start a statistics project. FinCompass is designed around pooled model families rather than requiring a separate model to be trained for every company. For the default supported US common-equity path, the intended flow is:
\begin{center}
\textbf{Enter ticker $\rightarrow$ identify market $\rightarrow$ choose benchmark $\rightarrow$ update missing data $\rightarrow$ choose the strongest applicable model $\rightarrow$ Forecast $\rightarrow$ Start Live.}
\end{center}
If the security is a supported US common stock, the software can use a pooled US-equity model even when that ticker was not one of the securities used to fit the model. The applicability checks still require the asset class, region, security type, benchmark family, horizon, and feature contract to match.

Guided mode should not ask the user to interpret phrases such as ``locked-test gate failed,'' ``feature contract,'' ``candidate rejected,'' or ``model ID.'' Those details belong in Research Details. If only a hard-valid Bayesian baseline is available, Guided mode should show the probability with \textbf{Limited evidence}. If a stronger validated model exists, it is preferred automatically. If fresh prices are missing, FinCompass should obtain or update the required local history before forecasting. A true unsupported asset is explained in plain language rather than being given a misleading probability.

\section{Forecast evidence labels}
Version 2.0 distinguishes model validity from stronger evidence of predictive skill.

\subsection{Limited evidence}
A \textbf{Bayesian baseline} model has enough valid data, uses only matured historical outcomes, fits numerically, and produces calibrated probabilities. It has not demonstrated enough out-of-sample improvement to qualify for the stronger research-validation tier.

You may see a Forecast such as:
\begin{quote}
54\% estimated probability of outperforming the S\&P 500 over the next year. Evidence: Limited.
\end{quote}
This is a usable model estimate, not a claim of proven alpha. Baseline models are not allowed to control adaptive Live adjustments.

\subsection{Validated research}
A validated-research model has passed the configured locked-test statistical gates. This is stronger evidence, but it can still have research-data limitations such as survivorship bias or incomplete corporate-action controls.

\subsection{Validated market}
This is the strongest current FinCompass tier. It requires the statistical gate plus stronger historical-universe and point-in-time data controls.

\section{Why a model may exist even when it does not ``beat the market''}
A probability model can be mathematically valid without showing strong incremental predictive skill. For example, a model can converge, remain leakage-free, and be reasonably calibrated while producing nearly the same proper score as the unconditional historical event rate. FinCompass labels such a model Limited evidence instead of calling the software broken or pretending the model is validated.

The stronger validation thresholds are not lowered. They determine the strength of the claim, not whether every coherent Bayesian estimate must disappear.

\section{Supported model horizons in the bundled research package}
Version 2.0 includes Bayesian reference artifacts for 6, 12, 24, and 36 months within the declared US common-equity / S\&P 500 domain. The 12-month enhanced reference model also carries a validated-research result on the bundled historical research corpus.

The presence of an artifact does not make every ticker scientifically supported. The ticker must match the model's declared market, asset class, security type, benchmark family, and feature contract.

\section{Unsupported securities}
FinCompass may still calculate ordinary analytics for an unsupported asset while refusing to produce a Forecast. For example, a bond ETF can have volatility and drawdown calculations even when no bond-forecast model has been validated.

This is intentional. A US-equity model is not automatically appropriate for Canada, Japan, ETFs, bonds, commodities, or cryptoassets.

\section{Data freshness}
Data freshness asks: \emph{How current is the local market history?}

A normal update begins shortly before the most recent local observation, downloads only the missing tail, compares the overlap, records provider corrections, and appends new observations. If the price basis changes---for example from raw prices to adjusted prices---FinCompass rebuilds that affected symbol rather than silently mixing incompatible histories.

\section{Model freshness}
Model freshness asks: \emph{How current was the information used when the model was trained?}

A model trained through 2022 can still load in 2026, but it does not contain the intervening market behavior. FinCompass therefore compares the training cutoff with current compatible data and can report states such as:
\begin{itemize}
\item current;
\item new data available;
\item retraining recommended;
\item stale;
\item validation refresh due.
\end{itemize}

Data freshness and model freshness are different.

\section{Updating a model}
An update should use the entire valid accumulated research corpus, not just the newly downloaded rows. The process is:
\begin{enumerate}
\item update missing data;
\item reconcile overlap and corrections;
\item rebuild the chronological dataset;
\item include only outcomes whose forward endpoint has matured;
\item train a new immutable candidate;
\item evaluate the candidate;
\item leave the current active model unchanged during this process.
\end{enumerate}

If a newer candidate is appropriate, FinCompass should offer \textbf{Use newer model} or \textbf{Keep current model}. Validation alone must never silently change the active model.

\section{Why future labels must mature}
Suppose the Forecast horizon is one year. A row observed today cannot have a one-year outcome tomorrow. The endpoint must actually occur. This simple rule prevents a major source of leakage.

New observations can be used to construct today's backward-looking features before they are eligible as historical supervised labels.

\section{Live}
Live begins from a frozen Forecast anchor. New information may produce an adaptive candidate, but the adaptive probability is applied only after enough completed outcomes, unique dates, temporal span, calibration evidence, fresh sources, and acceptable loss behavior exist.

A Limited-evidence Bayesian baseline can be tracked for freshness but does not control adaptive Live in Version 2.0.

\section{Performance analytics}
FinCompass includes deterministic calculations such as annualized return, volatility, Sharpe ratio, Sortino ratio, maximum drawdown, beta, and tracking error.

These are descriptions of historical behavior. They are not Forecast probabilities.

\section{Risk analytics}
Risk tools include historical and Gaussian Value at Risk, historical expected shortfall, EWMA volatility, and drawdown. VaR is a threshold under a stated method and confidence level; it is not the worst possible loss.

\section{Technical analytics}
Technical calculations include moving averages, RSI, MACD, Bollinger Bands, ATR, and related trend or volatility summaries. A technical indicator being available does not mean it is automatically used by the Forecast model.

\section{Valuation}
Discounted cash flow estimates present value under explicit assumptions. A simplified FCFF valuation is
\begin{equation}
V=\sum_{t=1}^{N}\frac{FCFF_t}{(1+WACC)^t}+\frac{TV_N}{(1+WACC)^N}.
\end{equation}
Because WACC, growth, margins, and terminal assumptions can dominate the output, use sensitivity analysis rather than treating one DCF number as truth.

Valuation and Forecast answer different questions.

\section{Fixed income}
The deterministic fixed-income module includes bond price, current yield, yield to maturity, Macaulay duration, modified duration, convexity, and related calculations from supplied inputs. Duration is an approximation of rate sensitivity; it is not a guarantee of realized price movement.

\section{Options}
The options module includes Black--Scholes--Merton pricing and major Greeks such as Delta, Gamma, Vega, Theta, and Rho. These are model-based theoretical quantities. Actual execution prices depend on spreads, liquidity, volatility surfaces, exercise style, dividends, and other market conditions.

\section{Portfolio analytics}
Portfolio calculations can combine holdings into weights, historical return, volatility, covariance/correlation, and concentration summaries. Individual Forecast probabilities must not be naively averaged and called a portfolio probability.

\section{Research Details: Bayesian model}
For feature vector $x$,
\begin{equation}
P(Y=1\mid x)=\frac{1}{1+\exp[-(\beta_0+x^\top\beta)]}.
\end{equation}
Shrinkage priors regularize the coefficients. Around the posterior mode, FinCompass uses a Laplace approximation to estimate coefficient uncertainty. Posterior draws are propagated to a probability interval and a separate validation block is used for calibration.


\section{Research Details: regime-aware Bayesian model}
Markets do not behave identically in every environment. Version 2.0 therefore also contains an experimental regime-aware Bayesian reference family. A three-state Gaussian hidden Markov model is fitted only on backward-looking benchmark features: one-month benchmark return, six-month benchmark return, and six-month benchmark volatility. It estimates filtered probabilities for three ordered market states that are described as risk-off, neutral, and risk-on.

The regime probabilities are then added to the small Bayesian reference feature set. The forecast remains a regularized Bayesian probability model; the hidden-state layer is not treated as a guaranteed trading signal. The HMM is fitted on training dates only, and later validation/test state probabilities are obtained by forward filtering so future observations cannot change earlier regime assignments.

This model family is packaged primarily as a Research-mode alternative. On the current historical corpus it did not consistently improve the stronger reference architecture, so it is not automatically preferred in Guided mode. This is an important design rule: adding model complexity does not earn promotion by itself.

\begin{center}
\begin{tabular}{rrrrr}\toprule
Horizon & Test $n$ & AUC & Brier skill & ECE\\\midrule
6 months & 539 & 0.444 & -0.045 & 0.086\\
12 months & 532 & 0.508 & -0.059 & 0.067\\
24 months & 518 & 0.485 & 0.009 & 0.029\\
36 months & 497 & 0.321 & -0.043 & 0.163\\
\bottomrule\end{tabular}
\end{center}
The 24-month regime model shows slightly positive proper-score skill in this locked sample, but this isolated result is not sufficient for promotion to a stronger evidence tier.

\section{Research Details: proper scores}
The Brier score for one binary outcome is
\begin{equation}
BS=(p-y)^2.
\end{equation}
Lower is better. Brier skill compares the model with a reference forecast. Log loss penalizes highly confident errors more strongly. ROC AUC describes ranking, not probability calibration. Expected calibration error describes the gap between predicted probabilities and observed frequencies within bins.

A good probability system should not rely on AUC alone.

\section{Research Details: locked test}
The locked test is useful only if it remains locked. Repeatedly trying feature sets and selecting the one that performs best on the locked test turns that test into another tuning dataset. FinCompass therefore separates training, validation/calibration, and locked-test evaluation.

\section{Research Details: the current bundled evidence}
The 12-month enhanced research model uses eight historical US equities against the S\&P 500. On 532 locked-test observations, the recorded metrics are approximately:
\begin{center}
\begin{tabular}{lr}\toprule
ROC AUC & 0.576\\
Brier skill & 0.036\\
Log-loss skill & 0.027\\
ECE & 0.045\\
Calibration slope & 1.27\\
\bottomrule\end{tabular}
\end{center}
The result passes the configured research gate, but the effect is modest and the dataset has important historical-universe limitations.

The 6-, 12-, 24-, and 36-month Bayesian reference artifacts are retained as Limited-evidence baselines. Their locked-test results do not support stronger predictive claims.

\section{Research Details: applicability}
Every forecast model should state:
\begin{itemize}
\item asset class;
\item region/market;
\item security type;
\item benchmark family;
\item feature frequency;
\item forecast horizon;
\item training universe and period.
\end{itemize}
Computing the required features is not enough. Scientific applicability must also match.

\section{Research Details: analytics versus features}
The analytics library is intentionally broader than the Forecast feature set. A new metric becomes a forecasting feature only after it receives a versioned definition, point-in-time review, missing-data policy, training experiment, and proper validation. This protects against accidental feature proliferation and data-snooping.

\section{What FinCompass does not claim}
FinCompass does not claim guaranteed returns, guaranteed outperformance, causal factor exposures, infallible DCF values, worst-case-loss bounds from VaR, or universal model validity across markets. It is a research and education tool designed to make both calculations and limitations inspectable.

\section{A practical checklist}
Before using a Forecast, ask:
\begin{enumerate}
\item What exact event is being forecast?
\item What is the horizon?
\item What benchmark is used?
\item Does the ticker match the model's applicability domain?
\item Is local data current?
\item Is the model current?
\item Is the evidence Limited, validated research, or validated market?
\item Is Live actually eligible, or is the frozen Forecast still controlling the display?
\end{enumerate}

\end{document}
